\sscp{Distribution of reflexes associated with **mantəR}{14-04-02}
The distribution of reflexes of **mantəR ‘cuscus’
and formally close words have significant methodological implications.
Very simply, cognates meaning ‘cuscus’ are found in limited parts
of the biogeographical regions where marsupials are naturally
found or have been introduced by humans prior to the arrival
of Austronesian-speaking peoples (see \frf{14-05}).
Some of these cognates are also in Sulawesi within the natural
range of several species of cuscus.
That fact presents a challenge for subgrouping.
If we apply the same logic as in \citet{Blust1982,
Blust1993,Blust2009,Blust2012b}, then we would be required to posit a subgroup
that includes EMP, languages of \tsc{Sula-Buru}, \tsc{\ili{Ambon}-Seram},
\ili{Banda}, \tsc{Seram-Tanimbar-Bomberai}, \tsc{Aru}, \tsc{Timor-Babar},
\tsc{\ili{Central} Timor}, excluding \tsc{\ili{Bima}-Lembata},
but also including the \tsc{\ili{Muna}-Buton} languages of south\ili{east} Sulawesi.
Yet there is no known basis in the historical phonology or morphosyntax
to support such a subgroup.
The problems with **mans(aə)r reflect some of the significant
drawbacks of basing a subgroup on lexical evidence.\footnote{
		Looking beyond the range of marsupials at wider possibilities,
		\citet[143, endnote 72]{Collins1983} linked **mans(aə)r ‘cuscus’ to
		a form further back up the chain,
		saying “it is almost certain that this form is related to certain Formosan forms.
		Tsuchida (personal communication) reconstructs *mangse(lR) pygmy
		deer based on such forms as \ili{Amis} \it{mangcɨl}.” We also note
		similar terms in western Indonesian languages associated with
		several varieties of the ‘pygmy deer’ or ‘barking deer’ (genus:
		\it{Muntiacus}), such as Sunda \it{mənʧək} ‘barking deer,
		pygmy deer, muncak deer’.
		But these forms probably ultimately trace back to \ili{Sanskrit} \it{mañcaka}
		‘deer’, including Javanese \it{menʤaŋan} ‘deer,
		kijang’.} 
\ili{As} can be seen in \frf{14-05},
the distribution of reflexes of **mantəR have very little overlap
with the reflexes of \it{k}-initial words for ‘cuscus’ shown in \frf{14-03}.
The \it{k}-initial family of forms are mostly found in Sulawesi,
Halmahera, parts of the Bird's Head,
and the \tsc{Eastern Islands} node of \tsc{Seram-Tanimbar-Bomberai},
while words related to **mantəR in \frf{14-05} are concentrated
in the central and southern parts of our target region,
on Buton Island (south\ili{east} Sulawesi),
and in some Papuan languages in Indonesian Papua.
(White circles or squares reflect words glossed ‘cuscus’ from
etyma other than **mantəR.) Note that there are no reflexes of
**mantəR in the \tsc{Eastern Islands} node of \tsc{Seram-Tanimbar-Bomberai},
since those languages relate to **kVndoR(a) (see \srf{14-03-04}).

\begin{figure}[h]\centering
	\includegraphics[width=\columnwidth]{manteR.pdf}
	\fcp{Reflexes of **mantəR and their glosses}{14-05}
\end{figure}

Tables \ref{tab:14-07a}, \ref{tab:14-07b}, and \ref{tab:14-07c} provide a sample of the forms,
distribution and glosses according to the sources for reflexes of **mantəR
and formally close words in Austronesian languages in the target
region, extended to include South\ili{east} Sulawesi.
Final vowels in \tsc{\ili{Ambon}-Seram} languages are extra-metrical
epenthetic vowels or noun markers.
In \crf{ch:STB}, we noted that in \tsc{Seram-Tanimbar-Bomberai}
languages, ultimate schwa *ə > \it{a}.
More complete data, including words that do not resemble **mantəR,
are found in appendix \ref{ch:SurMar}.
There are no reflexes of **mantəR ‘cuscus’ in the \tsc{Eastern
Islands} or \tsc{East Seram} nodes of \tsc{Seram-Tanimbar-Bomberai},
though data for \tsc{East Seram} is limited.
%Furthermore, there are no reflexes of **mantəR ‘cuscus’ in the
%\tsc{East Seram} node of \tsc{Seram-Tanimbar-Bomberai},
%although the data for that node are limited.

\begin{table}[p]\centering\tcp{Selected reflexes of **mantəR and similar words (part 1)}{14-07a}
\begin{adjustbox}{totalheight=\textheight}
	%\begin{threeparttable}
	\st{2pt}
	\wg\begin{tabularx}{\linewidth}{lll>{\raggedright\arraybackslash}X}\lsptoprule
Language	&	**mantəR	&	local gl.	&	Comment	\\	\midrule\hhline{~}
\mc{4}{c}{\tsc{\ili{Muna}-Buton} (\tsc{Celebic})}							\\	\midrule
\ili{Cia-Cia}	&	mante	&	cuscus	&	\cite[33]{MeadTruong2021}	\\
\ili{Kaisabu}	&	mante	&	bear cuscus	&	\cite[16]{Mead2017a}	\\	
\ili{Kaimbulawa}	&	mante	&	bear cuscus	&	David Mead, pers. comm.; also ‘dwarf cuscus’	\\	
\ili{Pancana}	&	mante	&	bear cuscus	&	David Mead, pers. comm.; also ‘dwarf cuscus’	\\	
\ili{Kioko}	&	mante	&	bear cuscus	&	\cite[13]{Mead2020c}	\\	
\ili{Kambowa}	&	mante	&	bear cuscus	&	\cite[14]{Mead2020b}	\\	
\ili{Busoa}	&	mante	&	cuscus	&	\cite[335]{vandenBerg2020}	\\	\midrule
\mc{4}{c}{\tsc{Sula-Buru}}							\\	\midrule
%\ili{Lisela}	&	\cg{(blaɸen)}	&	cuscus	&	\cite{GrimesGrimes2020}; loan?\su{†}	\\	
%\ili{Buru} 	&	\cg{(tonal)}	&	cuscus	&	\cite{GrimesGrimes2020}; loan\su{‡}	\\	
\ili{Ambelau}	&	mate	&	marsupial	&	\cite[36]{Collins1981} (`marsupial' = `cuscus')	\\	\midrule
\mc{4}{c}{\tsc{Seram-Tanimbar-Bomberai}}							\\	\midrule
\ili{Masiwang}	&	\cg{(mani)}	&	cuscus	&	\cite[358]{LeCocqdArmandville1901} from *manuk ‘bird’	\\	
\ili{Bati}	&	\cg{(manuk faniki)}	&	cuscus	&	\cite[196]{Ellen1993} from *manuk ‘bird’ and *paniki `bat'	\\	
\ili{Kei} Besar	&	mendar	&	cuscus	&	\cite[196]{Ellen1993}	\\	
\ili{Teor}	&	medar	&	cuscus	&	Craig Marshall (pers. comm. September 2025) \\%\cite{Travis-Grimes1990}	\\	
\ili{Fordata}	&	medar	&	cuscus	&	\cite{MarshallMarshall-Grimes1990}	\\	
\ili{Erokwanas}	&	matafeis	&	cuscus	&	\cite{Donohue2010}	\\	
\ili{Yamdena}	&	manr-e	&	(white) cuscus	&	\cite{MettlerMettler1996} ‘kuskus, putih’, \cite[65]{Drabbe1932b} ‘boombeer’ = ‘tree bear’	\\	\midrule
\mc{4}{c}{\tsc{Aru}}							\\	\midrule
\ili{Ujir}	&	meday	&	cuscus	&	\cite[263]{Schapper2011}	\\	
Goda-Goda	&	medah	&	cuscus	&	\cite[264]{Schapper2011}; village in West \ili{Kola} area	\\	
East \ili{Kola}	&	moduh,	&	cuscus	&	\cite{Nivens2017},	\\	\hhline{~}
	&	modah	&		&	\cite[264]{Schapper2011}	\\	
\ili{Ngaibor}	&	medar	&	marsupial	&	\cite[241]{Blust1982} village in SW \ili{Tarangan} area	\\	
SW \ili{Tarangan}	&	\cg{(man)}	&	bandicoot?	&	\cite{Nivens2017} question mark in original	\\
		\lspbottomrule
	\end{tabularx}
		%\begin{tablenotes}\tnsize
			%\item[†] {The sound correspondences are wrong for \ili{Lisela} \it{blaɸen}
								%‘cuscus’ to be from *balabaw ‘rat’ (expected form *[ɸlaɸa]).
								%It is possibly from Dutch \it{blafen} ‘bark,
								%bay, cry’, referring to the territorial night call/bark of the
								%male cuscus, often imitated by hunters near the full moon.}
			%\item[‡]	{The final l{\#} in \ili{Buru} \it{tonal} makes this form a suspect loan.
								%In Sulawesi, \tsc{\ili{Bungku}-Tolaki} languages \ili{Kulisusu}
								%and \ili{Mori Bawah}, and \tsc{\ili{Muna}-Buton} languages \ili{Cia-Cia},
								%\ili{Kaisabu}, \ili{Lasalimu}, \ili{Kumbewaha}, have \it{toŋali} ‘dwarf cuscus
								%(\it{Strigocuscus celebensis)}’ (David Mead pers. comm. Oct. 2020).
								%Some \tsc{\ili{Muna}-Buton} languages have *ŋ > \it{ŋ},
								%\it{n}, such as \ili{Kambowa} \it{tonali} ‘dwarf cuscus’.
								%There is a large
								%and long-standing presence of immigrants from South\ili{east} Sulawesi on \ili{Buru} \citep{Grimes2000a}.
								%\ili{Buru} tends to preserve \it{ŋ} in both retentions and loans.
								%So all in all,
								%\ili{Buru} \it{tonal} ‘cuscus’ appears to be from a south\ili{east} Sulawesi source.
								%This further evidence of a Sulawesi-Maluku link for borrowing
								%other terms for ‘cuscus’ in addition to **mantəR is significant.}
			%\item[Δ]
			%\item[‖]
			%\item[¶]
		%\end{tablenotes}
	%\end{threeparttable}
\end{adjustbox}
\end{table}

\begin{table}[p]\centering\tcp{Selected reflexes of **mantəR and similar words (part 2)}{14-07b}
%\begin{adjustbox}{totalheight=\textheight}
	%\begin{threeparttable}
	\wg\st{4pt}\begin{tabularx}{\linewidth}{ll>{\hspace{-\tabcolsep}}l<{\hspace{-\tabcolsep}}>{\raggedright\arraybackslash}X}\lsptoprule
Language	&	**mantəR	&	local gl.	&	Comment	\\	\midrule \hhline{~}
\mc{4}{c}{\tsc{\ili{Ambon}-Seram}}							\\	\midrule
\ili{Kayeli}	&	maren	&	cuscus	&	\cite[40]{Stresemann1927}	\\	
\ili{Manipa}	&	marele	&	marsupial	&	\cite[68]{Collins1983} (`marsupial' = `cuscus')	\\
\ili{Piru}	&	marel	&	cuscus	&	\cite[304f]{vanEkris1864}	\\	
\ili{Batumerah}	&	marela	&	marsupial	&	\cite[86]{Collins1983} (`marsupial' = `cuscus')	\\	
\ili{Larike}	&	máridu	&	cuscus	&	\cite[77]{LaidigLaidig1991}	\\	
\ili{Alang}	&	marilu	&	marsupial	&	\cite[84]{Collins1983} (`marsupial' = `cuscus')	\\	
\ili{Asilulu}	&	marel	&	cuscus	&	\cite[70]{Collins2003b}	\\	
\ili{Hitu}	&	makel	&	cuscus	&	\citet[68]{Collins1983}, \citet[196]{Ellen1993} (*nt > \it{k})	\\	
\ili{Kamarian}	&	maker	&	cuscus	&	\cite[305]{vanEkris1864} (*nt > \it{k})	\\	
\ili{Haruku}	&	makel	&	cuscus	&	\cite[305]{vanEkris1864} (*nt > \it{k})	\\	
\ili{Saparua}	&	makel	&	cuscus	&	\cite[305]{vanEkris1864} (*nt > \it{k})	\\	
\ili{Latu}	&	makell-o	&	marsupial	&	\cite[68]{Collins1983} (`marsupial' = `cuscus')	\\	
\ili{Amahai}	&	makerr-o	&	cuscus	&	\cite[43]{Stresemann1918}; (*nt > \it{k})	\\	
\ili{Paulohi}	&	makel-e	&	cuscus	&	\citet[38]{Stresemann1918}, \citet[196]{Ellen1993}	\\	
North \ili{Alune}	&	marel	&	cuscus	&	\citet[196]{Ellen1993}, \citet{TaguchiTaguchi-Grimes1990}; (*nt > \it{r}; cf. *punti > \it{uri} ‘banana’)	\\	
West \ili{Alune}	&	maʧel	&	cuscus	&	\cite[305]{vanEkris1864} (*nt > \it{ʧ}; cf. *punti > \it{uʧi} ‘banana’)	\\	
\ili{Yalahatan}	&	marer-e	&	cuscus	&	\cite{BeckleyND}	\\	
S. \ili{Wemale}	&	madel-e	&	marsupial	&	\cite[55]{Collins1983} (`marsupial' = `cuscus')	\\	
N. \ili{Wemale}	&	marel-e	&	marsupial	&	\cite[55]{Collins1983} (`marsupial' = `cuscus')	\\	
S. \ili{Nuaulu}	&	maran-e	&	cuscus	&	\citet[196]{Ellen1993}[65], \citet{MatokeBolton2005} *R/*r > \it{n}	\\	
%\ili{Manusela}	&	(ihisi)	&	cuscus	&	\cite{MacDonald2015}	\\	
\ili{Hoti}	&	mala	&	cuscus	&	\cite[256]{Ludeking1868}	\\	
\ili{Liana-Seti}	&	mair-a	&	marsupial	&	\citet{McCollum-Grimes1990}; metathesised from **mari	\\ \midrule\hhline{~}
\mc{4}{c}{Banda}							\\	\midrule
\ili{Banda}	&	mender,	&	cuscus	&	\cite[196f]{Ellen1993}	\\\hhline{~}
	&	mbender	&		&	\cite[45]{Kakerissa1996}	\\
		\lspbottomrule
	\end{tabularx}
		%\begin{tablenotes}\tnsize
			%\item[†] {}
		%\end{tablenotes}
	%\end{threeparttable}
%\end{adjustbox}
\end{table}

\begin{table}[p]\centering\tcp{Selected reflexes of **mantəR and similar words (part 3)}{14-07c}
%\begin{adjustbox}{width=\textwidth}
	\begin{threeparttable}\st{4pt}
	\wg\begin{tabularx}{\linewidth}{lll>{\raggedright\arraybackslash}X}\lsptoprule
Language	&	**mantəR	&	local gl.	&	Comment	\\	\midrule\hhline{~}
\mc{4}{c}{\tsc{Timor-Babar}}							\\	\midrule
\ili{Selaru}	&	mat	&	cuscus	&	\cite{CowardND}	\\	
\ili{Makatian}	&	mad-e	&	cuscus	&	Edwards fieldnotes	\\	
\ili{Teun}	&	mata	&	cuscus	&	Edwards fieldnotes	\\	
East \ili{Damer}	&	madar	&	cuscus	&	Edwards fieldnotes	\\	
SE. \ili{Babar}	&	mɔt	&	cuscus	&	Edwards fieldnotes	\\	
%E. \ili{Marsela}	&	mot	&	cuscus	&	Edwards fieldnotes	\\	
\ili{Emplawas}	&	mutt-o	&	cuscus, civet\su{†}	&	Edwards fieldnotes	\\	
Tela	&	mut-ɛ	&	cuscus	&	Edwards fieldnotes	\\	
\ili{Leti}	&	mada	&	cuscus	&	\citet[263]{Schapper2011}, \citet[423]{vanEngelenhoven2004}	\\	
\ili{Moa}	&	mada	&	cuscus	&	\cite[142]{Jonker1932} ‘buidelrat’	\\	
\ili{Roma}	&	madar,	&	cuscus	&	\cite[28]{Steven1991} (\it{matta} has	\\	\hhline{~}
	&	matta	&		&	  final VC → CV metathesis	\\	\hhline{~}
	&				&		&	  with C mutation/assimilation)	\\	
\ili{Kisar}	&	masla	&	cuscus	&	\cite{ChristensenChristensen-Grimes1990}	\\	
\ili{Makuva}	&	mate	&	opossum	&	\cite[433]{vanEngelenhoven2009}	\\	
Waima{\Q}a	&	mide	&	cuscus	&	\cite{HimmelmannFreitas2006}\su{‡}	\\	
Ili{\Q}uun	&	maʧa	&	cuscus	&	\cite[127]{DeJosselinDeJong1947}	\\	
\ili{Tugun}	&	maʧar	&	cuscus	&	\cite{Hinton2023}	\\	
Dadu{\Q}a	&	madah	&	cuscus	&	\cite[99]{Penn2006}	\\	
\ili{Galolen}	&	madar	&	cuscus	&	\cite[264]{Schapper2011}	\\	
\ili{Tetun}	&	meda	&	cuscus	&	\cite[264]{Schapper2011}	\\	
\ili{Ketun} (\tsc{Meto})	&	ʔmakuʔ,	&	cuscus	&	\cite[214]{Edwards2021}	\\	\hhline{~}
	&	ʔmaukaʔ	&		&	*nt > \it{k}; *punti > \it{uki}	\\	
Fatule{\Q}u (\tsc{Meto})	&	maukuʔ,	&	cuscus	&	\cite[214]{Edwards2021}	\\\hhline{~}
	&	maukaʔ	&		&	\hp{\,}	\\	\midrule\hhline{~}
%Roi{\Q}is \ili{Amarasi}	&	\cg{(urum)}	&	cuscus	&	\cite[92]{Edwards2021}	\\	
%Amfo{\Q}an (\ili{Meto})	&	\cg{(ukum)}	&	cuscus	&	\cite[92]{Edwards2021}	\\	
\mc{4}{c}{\tsc{\ili{Central} Timor}}		\\	\midrule
\ili{Welaun}	&	lamara	&	cuscus	&	\cite[52]{Edwards2019}	\\	
\ili{Kemak}	&	mada	&	cuscus	&	\cite[264]{Schapper2011}; Edwards fieldnotes	\\	
\ili{Mambae}	&	mat	&	cuscus	&	\cite[264]{Schapper2011} (possibly /maat/ met. from **mata)	\\	
		\lspbottomrule
	\end{tabularx}
		\begin{tablenotes}\tnsize
			\item[†]	{\ili{Emplawas} \it{mutt-o} `cuscus, civet' is probably semantic
								expansion due to language attrition.}
			\item[‡]	{\citet[264]{Schapper2011} gives Waima{\Q}a \it{meda} for ‘cuscus’.
								This is probably a loan from \ili{Tetun}.}
		\end{tablenotes}
	\end{threeparttable}
%\end{adjustbox}
\end{table}

With the exception of Southwest \ili{Tarangan} \it{man} ‘bandicoot?’,\footnote{
		Southwest \ili{Tarangan} \it{man} is given as ‘bandicoot?’ / ‘tikus tanah?’
		(question marks in the original).
		This word is included under the entry for *koyi which has a note
		at the beginning ``Check species of all data in this record!!!''
		\it{Tikus tanah} is the term for `bandicoot' in the local variety of \ili{Malay}.
		The exact identification may indeed be ‘bandicoot’,
		but clearly more investigation is needed.
		Several other species of small marsupial from the family Dasyuridae
		(quolls and dunnarts) occur in the \ili{Aru} Islands (\srf{14-01}).}
all of the other forms associated with **mantəR in our target region
and in SE Sulawesi refer to some type of ‘cuscus’,
with four as generic ‘marsupial’ (one of which is outside the
range of bandicoots, and two of which are from \citet{Collins1983} who
regularly glosses words for cuscus as ‘marsupial’).

Southwest \ili{Tarangan} \it{man} is also irregular compared to reflexes
of Proto-\ili{Aru} **medaR ‘cuscus’ in other \tsc{Aru} languages \citep[264]{Schapper2011}.
It is the only monosyllable,
the only form with a nasal,
the only form not reflecting *aR,
and the only \ili{Aru} form with the vowel \it{a} in the first syllable.
Furthermore, the normal reflex of *nt and *nd in Southwest \ili{Tarangan} is \it{d}.
Examples include: *paŋdan > **ɸandan > \it{odan} `pandanus' and
*bəntas `hack a passage through vegetation' > **-peda(s) > \it{-ɸeda} `lead'.
The primary source (\citealt{Nivens2017},
under headword **koyi) explicitly notes that the species of this entry needs checking.
Finally, of the nine \ili{Aru} languages for which we have words
for ‘bandicoot’, none are similar in form to Southwest \ili{Tarangan} \it{man} (see
\srf{14-05}).
All these factors show that Southwest \ili{Tarangan} \it{man} is not cognate.
It is a chance resemblance.
To summarise:

\begin{enumerate}
	\item	CMP has little or no support as a subgroup,
	\item	there is little archaeological or linguistic evidence to
				support Blust’s scenario of an Austronesian movement from Halmahera
				moving south and then west through Maluku to the Lesser Sundas,
	\item	the range of bandicoot habitat is limited to the New Guinea
				region and nearby islands,
	\item	the overwhelming consistency is the
				gloss ‘cuscus’
				%with only one irregular form in only one language
				%glossed ‘bandicoot’
				outside of Oceanic languages, and
	\item	there are formally close reflexes in several
				languages of South\ili{east} Sulawesi referring to ‘bear cuscus (\it{Ailurops
				ursinus})’ that retain a medial *NC cluster.
\end{enumerate}

Given that the associated cognate in Sulawesi forms
and glosses that are within the natural range of marsupials but
outside the CEMP area mean that this form is not an exclusively
shared innovation restricted to CEMP.
We must also support \citet{Schapper2011} that ‘cuscus’ is a far
more appropriate gloss than ‘bandicoot’ for **mans(aə)r outside
the EMP region.
Therefore putative **mans(aə)r `bandicoot' does
not provide subgrouping support for
the CEMP hypothesis.